\bookchapter{Παράρτη$\mu$α E: Ευχαριστίες και διευκρινίσεις}

Σε ευχαριστού$\mu$ε που κράτησες $\mu$έχρι τέλους αυτό το «$\mu$άθη$\mu$α κινητών»! Κατά τη συγγραφή, καταφύγα$\mu$ε σε πληθώρα επίση$\mu$ων εγχειριδίων κατασκευαστών, οδηγών χρήσης και δη$\mu$όσιων δη$\mu$οφιλών άρθρων τεχνολογίας, για να εξασφαλίσου$\mu$ε την ακρίβεια των πληροφοριών.

Η τεχνολογία εξελίσσεται $\mu$ε ραγδαίους ρυθ$\mu$ούς· το περιεχό$\mu$ενο του βιβλίου βασίζεται στις κυρίαρχες τεχνολογίες και προϊόντα της εποχής της συγγραφής. Όταν θα το διαβάζεις, ίσως έχουν ήδη ε$\mu$φανιστεί νέες λειτουργίες και προϊόντα. Ελπίζου$\mu$ε οι ιδέες και οι $\mu$έθοδοι που διδάσκει αυτό το βιβλίο να σε βοηθήσουν να αντι$\mu$ετωπίζεις $\mu$ε ηρε$\mu$ία και τα επό$\mu$ενα ψηφιακά «παιχνίδια» του $\mu$έλλοντος.

\booksection{Βασικά ση$\mu$εία της ενη$\mu$έρωσης 1.6}

\begin{itemize}
  \item Επέκταση του «Αξία και διαχείριση ψηφιακών περιουσιακών στοιχείων» και προσθήκη κεφαλαίου «Διαχείριση εικόνων και συστή$\mu$ατος αρχείων» (ονο$\mu$ατοδοσία ώρα+γεγονός, ταξινό$\mu$ηση και οργάνωση, αυτό$\mu$ατος καθαρισ$\mu$ός 72 ωρών, απο$\mu$άκρυνση δεδο$\mu$ένων EXIF).
  \item Βελτίωση του συστή$\mu$ατος προστασίας «καθαρής/προστατευτικής λειτουργίας» και κατά των «ηλεκτρονικών παρασίτων».
  \item Προσθήκη των συχνών σχη$\mu$άτων απάτης της εποχής AI (face swap / $\mu$ί$\mu$ηση φωνής / προσποίηση αυθεντίας / ψεύτικες επενδύσεις) και ιδεών προστασίας για να τα ξεπεράσεις.
  \item Ξαναγράψι$\mu$ο της προστασίας κατά του doxxing («άνοιγ$\mu$α κουτιού») και των επιθέσεων κοινωνικής $\mu$ηχανικής.
  \item Επέκταση του κεφαλαίου προσβασι$\mu$ότητας (πλήρεις λειτουργίες VoiceOver/TalkBack, $\mu$εγέθυνση και φίλτρα χρω$\mu$άτων, περιγραφή εικόνων $\mu$ε AI, ζωντανοί υπότιτλοι, δίοδος κλήσεων RTT/TTY).
  \item Ξαναγραφήσι$\mu$ο του πλαισίου prompts $\mu$ε τη λογική «ρόλος-πλαίσιο-εργασία-περιορισ$\mu$οί», $\mu$ε 3 πλήρη παραδείγ$\mu$ατα σεναρίων.
  \item Προστέθηκαν τρία νέα προχωρη$\mu$ένα τεχνικά παράρτη$\mu$α: Παράρτη$\mu$α του Κεφαλαίου 1: ανάλυση της επέκτασης $\mu$νή$\mu$ης σε κινητά Android (ουσιαστικός ορισ$\mu$ός, αρχή λειτουργίας zRAM/swap, φράγ$\mu$α ταχύτητας, όρια εφαρ$\mu$ογής, ονο$\mu$ατοδοσία και ενεργοποίηση ανά κατασκευαστή).
  \item Παράρτη$\mu$α του Κεφαλαίου 20: Termux — φορητό περιβάλλον τερ$\mu$ατικού Linux σε συσκευές Android (βασικές λειτουργίες, οικοσύστη$\mu$α plugins, εγκατάσταση και ρύθ$\mu$ιση, τυπικά σενάρια όπως mobile development, διαχείριση διακο$\mu$ιστών, δοκι$\mu$ές ασφαλείας).
  \item Παράρτη$\mu$α του Κεφαλαίου 23: οδηγός εκκίνησης για το GitHub — από την αποθήκευση κώδικα στην πλατφόρ$\mu$α συνεργασίας (διάκριση Git vs GitHub, βασικές έννοιες, εγγραφή και πρώτα βή$\mu$ατα, ροή συνεργασίας $\mu$ε Pull Request, υπενθυ$\mu$ίσεις ασφαλείας, λύσεις πρόσβασης από την Κίνα).
\end{itemize}

\booksection{Βασικά ση$\mu$εία της ενη$\mu$έρωσης 1.7}

1. Προστέθηκαν 4 νέα κεφάλαια

\begin{itemize}
  \item Κεφάλαιο 29: Ψηφιακή υγεία — ανάλυση χρόνου χρήσης οθόνης, όρια εφαρ$\mu$ογών, λειτουργία «Μην ενοχλείτε», κανόνας 20-20-20 για την όραση, αυτοσωτηρία του αυχένα, διαχείριση ύπνου, «εκλεκτική τύφλωση» από άγχος πληροφοριών, πρόκληση ψηφιακής υγείας $\mu$ίας εβδο$\mu$άδας.
  \item Κεφάλαιο 30: Ασύρ$\mu$ατη προβολή (cast) και συστή$\mu$ατα αυτοκινήτου — σύγκριση των τριών «στρατοπέδων» AirPlay / Miracast / DLNA, ενσύρ$\mu$ατες λύσεις προβολής, λειτουργία επιφάνειας εργασίας Samsung DeX, οι τέσσερις $\mu$εγάλες λύσεις αυτοκινήτου CarPlay / Android Auto / HiCar / CarLife, πρακτική χρήση προβολής σε αίθουσα συσκέψεων.
  \item Κεφάλαιο 31: Δη$\mu$ιουργία βίντεο $\mu$ικρής διάρκειας — οι τρεις ερωτήσεις για την επιλογή θέ$\mu$ατος, ρύθ$\mu$ιση γυρίσ$\mu$ατος σε 5 βή$\mu$ατα, πέντε κινήσεις κά$\mu$ερας (προσέγγιση, απο$\mu$άκρυνση, πάν, $\mu$ετατόπιση, παρακολούθηση), ροή εργασίας έξι βη$\mu$άτων στο CapCut, βασικά ση$\mu$εία κατασκευής εξωφύλλου, η φιλοσοφία περιεχο$\mu$ένου «οι πρώτες 3 δευτερόλεπτα κρίνουν την τύχη».
  \item Κεφάλαιο 32: Εγχειρίδιο πρώτων βοηθειών ασφαλείας — SOP των «χρυσών» 30 λεπτών $\mu$ετά την απώλεια, διαδικασία 72 ωρών $\mu$ετά την απάτη, έλεγχος σε safe mode $\mu$ετά από $\mu$όλυνση, έξι βή$\mu$ατα ανάκτησης παραβιασ$\mu$ένου λογαριασ$\mu$ού, αντι$\mu$ετώπιση απειλητικών $\mu$ηνύ$\mu$ατα εκβίασης, πέντε καθη$\mu$ερινά $\mu$έτρα πρόληψης.
\end{itemize}

2. Προστέθηκαν δύο νέα κεφάλαια

\begin{itemize}
  \item Κεφάλαιο 27: Φιλικότητα προς ηλικιω$\mu$ένους και οικογενειακή φροντίδα — πίνακας σύγκρισης «απλής λειτουργίας» ανά $\mu$άρκα, οι έξι βασικές απάτες και οι τρεις γρα$\mu$$\mu$ές ά$\mu$υνας, απο$\mu$ακρυσ$\mu$ένη βοήθεια (ChangLian / Sunflower), λειτουργία φροντίδας του WeChat και φωνητική πληκτρολόγηση, SOS / ανίχνευση πτώσης / υπενθυ$\mu$ίσεις φαρ$\mu$άκων.
  \item Κεφάλαιο 28: Οδηγός αγοράς — $\mu$έθοδος επιλογής κινητού σε 3 βή$\mu$ατα, πίνακας χαρακτηριστικών «$\mu$όνο 5 στοιχεία», προτάσεις σε 5 επίπεδα προϋπολογισ$\mu$ού, πίνακας επιλογής για 7 σενάρια χρήσης, οι 5 $\mu$εγάλες παγίδες στην αγορά κινητού, ανάλυση της κατάλληλης στιγ$\mu$ής και των καναλιών αγοράς.
\end{itemize}

3. Νέες λειτουργίες ανάγνωσης και αλληλεπίδρασης

\begin{itemize}
  \item Πλωτή γρα$\mu$$\mu$ή ελέγχου ανάγνωσης: υποστηρίζει αναπαραγωγή/παύση, προηγού$\mu$ενο/επό$\mu$ενο τ$\mu$ή$\mu$α, ρύθ$\mu$ιση ταχύτητας 0.5x-2.5x, διακοπή.
  \item Γρήγορο κου$\mu$πί ακουστικών δεξιά: εκκίνηση ανάγνωσης $\mu$ε ένα κλικ, αυτό$\mu$ατη έναρξη από την τρέχουσα θέση ανάγνωσης.
  \item Συντο$\mu$εύσεις πληκτρολογίου: R για εναλλαγή ανάγνωσης, ESC για διακοπή.
  \item Βελτίωση ε$\mu$πειρίας ανάγνωσης: αυτό$\mu$ατη συνεχό$\mu$ενη $\mu$ετάβαση στο επό$\mu$ενο κεφάλαιο και αυτό$\mu$ατη απόκρυψη των στοιχείων ελέγχου ανάγνωσης κατά την εκτύπωση.
\end{itemize}

4. Ενη$\mu$ερώσεις της δο$\mu$ής

Πίνακας περιεχο$\mu$ένων, διαδρο$\mu$ή $\mu$άθησης (4 βή$\mu$ατα5 βή$\mu$ατα), χάρτης κεφαλαίων, αριθ$\mu$ός έκδοσης (1.61.7), αρίθ$\mu$ηση παραρτη$\mu$άτων (VVI) και metadata ενη$\mu$ερώθηκαν όλα συγχρόνως.

5. Ενη$\mu$έρωση εγκυρότητας του περιεχο$\mu$ένου

\begin{itemize}
  \item Πρότυπο επεξεργαστών: Snapdragon 8 Elite Gen5 / Dimensity 9500 κατατάσσονται στην κορυφαία κατηγορία (βάση 2026).
  \item Αναφορά ναυαρχίδων: iPhone 17 Pro, Samsung S26 Ultra, Xiaomi 17 Ultra, Huawei Pura 80 Ultra, vivo X300 Pro.
  \item Κεφάλαια AI: ενη$\mu$έρωση στα DeepSeek R1, ChatGPT-5 και στο πρότυπο υπολογιστικής ισχύος NPU 80 TOPS σε επίπεδο συσκευής.
  \item Διπλώνεται η οθόνη (foldables): η ζώνη τι$\mu$ών κατέβηκε έως τα 4.000-5.000 γιουάν.
  \item Οδηγός αγοράς: όλες οι κατηγορίες προϋπολογισ$\mu$ού και οι προτάσεις σεναρίων ανανεώθηκαν στα δεδο$\mu$ένα της αγοράς του 2026.
  \item Δίκτυο και οικοσύστη$\mu$α: Wi-Fi 7 / 5G-A / οικοσύστη$\mu$α HarmonyOS / LE Audio ενη$\mu$ερώθηκαν όλα στα τελευταία πρότυπα.
\end{itemize}

6. Νέο λειτουργικό υποσύστη$\mu$α: πίνακας εκ$\mu$άθησης (κάτω από την περιοχή Hero)

\begin{itemize}
  \item Κυκλικό διάγρα$\mu$$\mu$α προόδου ανάγνωσης: ε$\mu$φανίζει σε πραγ$\mu$ατικό χρόνο το ποσοστό των κεφαλαίων που έχεις διαβάσει, $\mu$ε SVG κυκλική κίνηση.
  \item Αυτοέλεγχος υγείας κινητού: λίστα ελέγχου 8 ση$\mu$είων ασφαλείας (κωδικός, αυθεντικοποίηση δύο παραγόντων, εντοπισ$\mu$ός κινητού, backup στο cloud, backup WeChat, ασφάλεια πληρω$\mu$ών, app κατά της απάτης, καθαρισ$\mu$ός ρυθ$\mu$ίσεων πληρω$\mu$ής χωρίς κωδικό).
  \item Μονι$\mu$οποίηση κατάστασης: τα τικαρισ$\mu$ένα αποθηκεύονται στο localStorage, η βαθ$\mu$ολογία υγείας υπολογίζεται δυνα$\mu$ικά και ε$\mu$φανίζεται $\mu$ε ράβδο προόδου σε διαβάθ$\mu$ιση χρώ$\mu$ατος.
  \item Μηνύ$\mu$ατα ανά επίπεδο: ανάλογα $\mu$ε τη βαθ$\mu$ολογία ε$\mu$φανίζεται διαφορετική ανατροφοδότηση («πλήρης βαθ$\mu$ός/πολύ καλά/$\mu$έτρια/κίνδυνος ασφαλείας» κ.λπ.).
\end{itemize}

\booksection{Βασικά ση$\mu$εία της ενη$\mu$έρωσης 1.8}

Προστέθηκε το έβδο$\mu$ο $\mu$έρος «Πρακτικός Προγρα$\mu$$\mu$ατισ$\mu$ός $\mu$ε AI» (Κεφάλαια 37-47, 11 κεφάλαια συνολικά): επανάσταση παραδείγ$\mu$ατος και Agentic Engineering, οι πέντε νοητικές αρχές προγρα$\mu$$\mu$ατισ$\mu$ού $\mu$ε AI, $\mu$ηχανική πλαισίου (context) και ανάπτυξη $\mu$ε προδιαγραφές (SDD), $\mu$ηχανική απαιτήσεων, πανόρα$\mu$α εργαλείων προγρα$\mu$$\mu$ατισ$\mu$ού $\mu$ε AI του 2026, πανόρα$\mu$α $\mu$εγάλων $\mu$οντέλων και πρακτική $\mu$ε API, $\mu$ηχανική συνο$\mu$ιλίας, πλήρης ροή έργου από το 0 στο 1, εξασφάλιση ποιότητας, ο δρό$\mu$ος της προόδου (MCP/συνεργασία πολλαπλών Agent) και εγχειρίδιο ταχείας αναφοράς. Ενη$\mu$ερώθηκαν συγχρόνως ο πίνακας περιεχο$\mu$ένων, η διαδρο$\mu$ή $\mu$άθησης (5 βή$\mu$ατα6 βή$\mu$ατα), ο χάρτης κεφαλαίων και ο δείκτης ταχείας αναζήτησης. Το περιεχό$\mu$ενο προέρχεται από το «Σκέψη Προγρα$\mu$$\mu$ατισ$\mu$ού $\mu$ε AI: Μεθοδολογικός Οδηγός από το Μηδέν στην Πλήρη Πρακτική», έκδοση Αυγούστου 2026.

Προστέθηκαν το όγδοο $\mu$έρος «Πρακτική Ανάπτυξη Android» (Κεφάλαια 48-58, 11 κεφάλαια) και το ένατο $\mu$έρος «Πρακτική Ανάπτυξη iOS» (Κεφάλαια 59-69, 11 κεφάλαια): και οι δύο πλατφόρ$\mu$ες καλύπτουν στήσι$\mu$ο περιβάλλοντος (Android Studio / Xcode), ταχύρρυθ$\mu$η εκ$\mu$άθηση γλώσσας (Kotlin / Swift), frameworks διεπαφής (Jetpack Compose / SwiftUI), δεδο$\mu$ένα και δίκτυο, δύο πλήρη πρακτικά έργα (εφαρ$\mu$ογή εξόδων + αναζήτηση καιρού), δοκι$\mu$ές και ανέβασ$\mu$α σε κατάστη$\mu$α (Google Play / App Store) και εγχειρίδιο ταχείας αναφοράς. Ενη$\mu$ερώθηκαν συγχρόνως ο πίνακας περιεχο$\mu$ένων, η διαδρο$\mu$ή $\mu$άθησης (6 βή$\mu$ατα8 βή$\mu$ατα), ο χάρτης κεφαλαίων και ο δείκτης ταχείας αναζήτησης. Το περιεχό$\mu$ενο προέρχεται από τα «Ανάπτυξη Android από το Μηδέν» και «Ανάπτυξη iOS από το Μηδέν», έκδοση Αυγούστου 2026.

\booksection{Συ$\mu$πληρώσεις περιεχο$\mu$ένου: χρήσι$\mu$ες εφαρ$\mu$ογές και φιλτράρισ$\mu$α ενοχλητικών κλήσεων}

\begin{itemize}
  \item Στο Κεφάλαιο 13 προστέθηκε η ενότητα 6 «Χρήσι$\mu$α εργαλεία συστή$\mu$ατος»: εργαλεία καθαρισ$\mu$ού όπως το Qingzhuo, ο ενσω$\mu$ατω$\mu$ένος καθαρισ$\mu$ός αποθηκευτικού χώρου του συστή$\mu$ατος, η Εθνική Υπηρεσία Κατά της Απάτης, οι Phone Manager των $\mu$αρκών, η προστασία από ενοχλητικές κλήσεις των παρόχων και το AdGuard, $\mu$ε την αρχή επιλογής «πρώτα το σύστη$\mu$α / επίση$\mu$α για συ$\mu$πλήρωση / συγκράτηση $\mu$ε τα εργαλεία τρίτων».
  \item Πίνακας σύγκρισης ενσω$\mu$ατω$\mu$ένου φιλτραρίσ$\mu$ατος ενοχλητικών ανά $\mu$άρκα: καλύπτει iOS, Xiaomi/Redmi, Huawei/Honor, OPPO/realme/OnePlus, vivo/iQOO, Samsung — διαδρο$\mu$ές πρόσβασης και προτεινό$\mu$ενες ρυθ$\mu$ίσεις ενεργοποίησης.
  \item Επέκταση των δυνατοτήτων φιλτραρίσ$\mu$ατος ενοχλητικών στο Κεφάλαιο 8: προστέθηκαν πρακτικές οδηγίες για αναγνώριση κλήσεων, λίστες $\mu$αύρης/λευκής λίστας, φιλτράρισ$\mu$α SMS $\mu$ε λέξεις-κλειδιά, σίγαση αγνώστων στο iOS, σύνδεση $\mu$ε «Μην ενοχλείτε» κ.λπ., $\mu$ε διασύνδεση προς το Κεφάλαιο 13.
  \item Πρόταση τριπλής προστασίας: φιλτράρισ$\mu$α συστή$\mu$ατος + προστασία από παρόχο + Εθνική Υπηρεσία Κατά της Απάτης, $\mu$ε έ$\mu$φαση στη λευκή λίστα για να αποφεύγονται λάθος $\mu$πλοκαρίσ$\mu$ατα.
\end{itemize}

\booksection{Συ$\mu$πληρώσεις περιεχο$\mu$ένου: προειδοποιήσεις καιρού / φυσικών καταστροφών / σεισ$\mu$ών}

\begin{itemize}
  \item Στο Κεφάλαιο 13 προστέθηκε η ενότητα 7: προειδοποιήσεις καιρού (επίση$\mu$α $\mu$ετεωρολογικά κανάλια + app καιρού του συστή$\mu$ατος + χρω$\mu$ατικά επίπεδα), προειδοποιήσεις φυσικών καταστροφών (πολλαπλά κανάλια $\mu$ετεωρολογίας/έκτακτης ανάγκης), αρχή λειτουργίας και παρανοήσεις των σεισ$\mu$ικών προειδοποιήσεων (προειδοποίηση $\neq$ πρόγνωση).
  \item Πίνακας ενεργοποίησης σεισ$\mu$ικών προειδοποιήσεων ανά $\mu$άρκα: Xiaomi, Huawei/Honor, οικοσύστη$\mu$α OPPO, vivo, iPhone, Samsung — συνηθισ$\mu$ένες διαδρο$\mu$ές και ση$\mu$εία προσοχής για τα δι και ώ$\mu$ατα.
  \item Στο Κεφάλαιο 32 προστέθηκε το «Σενάριο $\mu$ηδέν»: σειρά ελιγ$\mu$ών προστασίας $\mu$ετά από σεισ$\mu$ικό συναγερ$\mu$ό, αντι$\mu$ετώπιση έντονης βροχής/τυφώνα, σύνδεση $\mu$ε δορυφορικό SOS / επικοινωνία χωρίς δίκτυο· το εγχειρίδιο πρώτων βοηθειών συνδέεται διαθε$\mu$ατικά $\mu$ε τις ρυθ$\mu$ίσεις προειδοποιήσεων του Κεφαλαίου 13.
  \item Διασύνδεση της υποενότητας καιρού: η εγγραφή του app καιρού στις ταξιδιωτικές $\mu$ετακινήσεις ε$\mu$πλουτίστηκε $\mu$ε τις ειδοποιήσεις προειδοποίησης και την υπενθύ$\mu$ιση ότι «θα τις λα$\mu$βάνεις $\mu$όνο αν έχεις ανοίξει τις ειδοποιήσεις».
\end{itemize}

\booksection{Συ$\mu$πληρώσεις περιεχο$\mu$ένου: συνεργασία $\mu$εταξύ συσκευών και εργαλεία αποδοτικότητας}

\begin{itemize}
  \item Ενότητα 8 του Κεφαλαίου 13: έλεγχος υπολογιστή από το κινητό (Sunflower, ToDesk) και οι «κόκκινες γρα$\mu$$\mu$ές» κατά της απάτης· γραφείο $\mu$ε AI (Trae / WorkBuddy / Dumate) και ροή εργασίας τηλεχειρισ$\mu$ού υπολογιστή από κινητό· τη δυνα$\mu$ική ταπετσαρία «Artificial Desktop» της miHoYo.
  \item Γρήγορη $\mu$εταφορά σε τοπικό δίκτυο και clipboard: LocalSend, IP Messenger, AirDrop· KDE Connect, καθολικό clipboard, $\mu$εταφορά κινητό-υπολογιστής ανά $\mu$άρκα.
  \item Δυνα$\mu$ικοί κωδικοί 2FA: Aegis, Microsoft Authenticator, backup codes και ση$\mu$εία προσοχής κατά την αλλαγή συσκευής.
  \item TickTick και σάρωση χάρτινων εγγράφων: ανάλυση της χρήσης για πρόγρα$\mu$$\mu$α/εργασίες· συ$\mu$πλήρωση για εργαλεία σάρωσης και τεχνικές φωτογράφισης.
\end{itemize}

Σου ευχό$\mu$αστε να κυριαρχείς στο κινητό σου, να απολα$\mu$βάνεις την άνεση και τη χαρά που φέρνει η τεχνολογία, και να γίνεις ένας ξύπνιος, σίγουρος ψηφιακός πολίτης.

— Ο οδηγός σου για την ψηφιακή ζωή
